\chapter{Bilangan Desimal}\label{ch:g5:decimals}

Persepuluhan dan perseratusan muncul pada \cref{ch:g4:decimals};
menambahkan satu tingkat lagi, perseribuan, melengkapi sistem
desimal. Bab ini memantapkan cara membaca, menguraikan dan ---
terutama --- membandingkan bilangan desimal tanpa terjatuh ke
perangkapnya yang terkenal.

\section{Turun sampai perseribuan}

\begin{definition}[Tempat desimal]\label{def:g5:decimals:places}
Sesudah \omterm{def:g4:decimals:point}{titik desimal} datang persepuluhan, perseratusan, lalu
perseribuan --- masing-masing bernilai sepuluh kali lebih kecil
daripada tempat sebelumnya:
\[
4.362 = 4 + \frac{3}{10} + \frac{6}{100} + \frac{2}{1000}
= \frac{4\,362}{1000}.
\]
\end{definition}

\begin{omfigure}
\begin{tikzpicture}[scale=0.9]
  \draw (0,0) grid[xstep=2.3, ystep=0.85] (9.2,1.7);
  \node[font=\small] at (1.15,1.275) {satuan};
  \node[font=\small] at (3.45,1.275) {persepuluhan};
  \node[font=\small] at (5.75,1.275) {perseratusan};
  \node[font=\small] at (8.05,1.275) {perseribuan};
  \node[font=\large, omDef] at (1.15,0.425) {$4$};
  \node[font=\large, omThm] at (3.45,0.425) {$3$};
  \node[font=\large, omThm] at (5.75,0.425) {$6$};
  \node[font=\large, omThm] at (8.05,0.425) {$2$};
  \fill (2.3,0.06) circle (2.5pt);
\end{tikzpicture}

{\small Tabel nilai tempat dari $4.362$: bagian bulat berwarna biru,
bagian desimal berwarna merah, titiknya di antara satuan dan persepuluhan.}
\end{omfigure}

\begin{example}[Banyak bacaan, satu bilangan]\label{ex:g5:decimals:readings}
$4.362$ dapat dibaca ``empat titik tiga enam dua'', atau ``empat dan
$362$ perseribu'', atau diuraikan menjadi $4 + 0.3 + 0.06 + 0.002$.
Menambahkan \omterm{def:g1:counting:zero}{nol} di ujung tidak mengubah apa pun: $4.362 = 4.3620$.
Tetapi $4.362 \neq 4.0362$: \omterm{def:g1:counting:zero}{nol} \emph{di dalam} mendorong setiap
angka ke tempat yang lebih kecil.
\end{example}

\begin{example}[Di mana letaknya pada garis?]\label{ex:g5:decimals:line}
Untuk meletakkan $4.362$: antara $4$ dan $5$; diperbesar, antara
$4.3$ dan $4.4$; diperbesar lagi, antara $4.36$ dan $4.37$, lebih
dekat ke $4.36$. Setiap angka desimal adalah satu tingkat perbesaran lagi.
\end{example}

\begin{omfigure}
\begin{tikzpicture}[scale=1.15]
  \draw[-Stealth] (-0.2,1.6) -- (10.4,1.6);
  \foreach \i in {0,...,10} \draw (\i,1.52) -- (\i,1.68);
  \node[below, font=\small] at (0,1.5) {$4$};
  \node[below, font=\small] at (10,1.5) {$5$};
  \node[below, font=\small] at (3,1.5) {$4.3$};
  \node[below, font=\small] at (4,1.5) {$4.4$};
  \fill[omThm] (3.62,1.6) circle (1.7pt);
  \draw[omExo, thin] (3,1.45) -- (0,0.15);
  \draw[omExo, thin] (4,1.45) -- (10,0.15);
  \draw[-Stealth] (-0.2,0) -- (10.4,0);
  \foreach \i in {0,...,10} \draw (\i,-0.08) -- (\i,0.08);
  \node[below, font=\small] at (0,-0.1) {$4.3$};
  \node[below, font=\small] at (10,-0.1) {$4.4$};
  \node[below, font=\small] at (6,-0.1) {$4.36$};
  \fill[omThm] (6.2,0) circle (1.7pt)
    node[above, font=\small] {$4.362$};
\end{tikzpicture}

{\small Memperbesar garis bilangan: rentang dari $4.3$ sampai $4.4$,
setelah diperbesar, ternyata terpotong lagi menjadi sepuluh
perseratusan --- dan $4.362$ terletak tepat setelah $4.36$.}
\end{omfigure}

\section{Membandingkan bilangan desimal}

\begin{method}[Membandingkan dua bilangan desimal]\label{met:g5:decimals:compare}
\begin{enumerate}
  \item Bandingkan bagian bulatnya: $6.1 > 5.987$.
  \item Bagian bulatnya sama: bandingkan bagian desimalnya angka
        demi angka mulai dari titiknya --- persepuluhan, lalu
        perseratusan, lalu perseribuan;
  \item menambahkan \omterm{def:g1:counting:zero}{nol} di ujung membuat perbandingannya adil:
        $2.5$ lawan $2.48$ menjadi $2.50$ lawan $2.48$, dan
        $50 > 48$ perseratusan.
\end{enumerate}
Perangkapnya, untuk yang terakhir kali: \emph{bagian desimal yang
lebih panjang tidak berarti bilangannya lebih besar} --- $2.48$ punya angka lebih banyak daripada $2.5$ dan justru lebih kecil.
\end{method}

\begin{example}\label{ex:g5:decimals:order}
Urutkan $7.3$; $7.09$; $7.31$; $7.299$. Lengkapi sampai tiga angka
desimal: $7.300$; $7.090$; $7.310$; $7.299$. Lalu
\[
7.09 < 7.299 < 7.3 < 7.31 .
\]
Perhatikan $7.299 < 7.3$ walaupun $299$ terlihat besar: $299$
perseribu melawan $300$ perseribu.
\end{example}

\begin{example}[Menyelipkan bilangan di antara dua desimal]\label{ex:g5:decimals:between}
Adakah bilangan antara $5.7$ dan $5.8$? Ada, banyak sekali: $5.75$,
$5.71$, $5.799$\,\dots\ Antara $5.79$ dan $5.8$? Banyak lagi:
$5.795$, misalnya. Di antara dua bilangan desimal yang berbeda, kita
\emph{selalu} dapat menyelipkan satu lagi --- tidak ada bilangan
desimal ``berikutnya''.
\end{example}

\section{Membulatkan bilangan desimal}

\begin{definition}[Pembulatan]\label{def:g5:decimals:round}
Untuk \omterm{def:g4:numbers:round}{membulatkan} ke satuan (atau ke persepuluhan, atau ke
perseratusan), lihatlah angka berikutnya: $5$ atau lebih \omterm{def:g4:numbers:round}{dibulatkan}
ke atas, selain itu \omterm{def:g4:numbers:round}{dibulatkan} ke bawah. Jadi $4.362$ membulat
menjadi $4$ (satuan), $4.4$ (persepuluhan), $4.36$ (perseratusan).
\end{definition}

\begin{example}[Uang dibulatkan ke sen]\label{ex:g5:decimals:money}
Harga yang terhitung $7.4983$ ditampilkan sebagai $7.50$: \omterm{def:g1:addition:def}{jumlah}
uang dalam kehidupan nyata \omterm{def:g4:numbers:round}{dibulatkan} ke perseratusan. Perhatikan
riaknya: $7.4983 \to
7.50$, ketika $49$ menjadi $50$ --- \omterm{def:g4:numbers:round}{pembulatan} dapat merambat
melewati beberapa angka.
\end{example}

\section{Latihan}

\begin{exercise}[$\star$]\label{exo:g5:decimals:1}
Tulislah sebagai bilangan desimal: $5 + \frac{2}{10} + \frac{7}{1000}$;
\quad $\frac{85}{100}$; \quad $\frac{4\,507}{1000}$; \quad ``dua
belas dan sembilan perseribu''.
\end{exercise}

\begin{exercise}[$\star$]\label{exo:g5:decimals:2}
Pada $27.418$: apa yang dibilang angka $4$? Angka $8$? Uraikan
bilangan itu seperti pada \cref{def:g5:decimals:places}.
\end{exercise}

\begin{exercise}[$\star$]\label{exo:g5:decimals:3}
Benar atau salah? Jelaskan dengan nilai tempatnya.
$3.60 = 3.6$; \quad $0.5 = 0.05$; \quad $2.070 = 2.07$; \quad
$1.3 = 1.03$.
\end{exercise}

\begin{exercise}[$\star$]\label{exo:g5:decimals:4}
Salin dan lengkapi dengan $<$, $>$ atau $=$:
\[
4.8 \;?\; 4.79, \qquad
0.302 \;?\; 0.32, \qquad
6.5 \;?\; 6.500, \qquad
9.09 \;?\; 9.1 .
\]
\end{exercise}

\begin{exercise}[$\star$]\label{exo:g5:decimals:5}
Urutkan dari yang terkecil ke yang terbesar: $3.14$; \ $3.4$; \ $3.041$; \
$3.104$; \ $3.41$.
\end{exercise}

\begin{exercise}[$\star$]\label{exo:g5:decimals:6}
Setiap bilangan berikut terletak di antara dua bilangan bulat yang
mana? Di antara dua bilangan satu angka desimal yang mana? \quad
$6.83$; \quad $0.492$; \quad $12.06$.
\end{exercise}

\begin{exercise}[$\star$]\label{exo:g5:decimals:7}
Bulatkan $8.276$: ke satuan; ke persepuluhan; ke perseratusan.
Bulatkan $3.95$ ke persepuluhan.
\end{exercise}

\begin{exercise}[$\star$]\label{exo:g5:decimals:8}
Berikan sebuah bilangan yang terletak tepat di antara: $2.6$ dan
$2.7$; \quad $0.41$ dan $0.42$; \quad $5.99$ dan $6$.
\end{exercise}

\begin{exercise}[$\star$]\label{exo:g5:decimals:9}
Tinggi empat tanaman adalah $0.85$ m, $1.2$ m, $0.9$ m dan $1.02$ m.
Urutkan dari yang terpendek ke yang tertinggi.
\end{exercise}

\begin{exercise}[$\star\star$]\label{exo:g5:decimals:10}
Lusi berkata: ``$7.12 > 7.9$ karena $12 > 9$''. Betulkan
kesalahannya, dengan membandingkan persepuluhannya lebih dulu.
\end{exercise}

\begin{exercise}[$\star\star$]\label{exo:g5:decimals:11}
Carilah semua bilangan dengan tepat dua angka desimal yang membulat
menjadi $6.4$ ketika \omterm{def:g4:numbers:round}{dibulatkan} ke persepuluhan. Berapa yang
terkecil dan berapa yang terbesar?

\end{exercise}
